\documentclass[12pt,a4paper]{article}

\usepackage[T1]{fontenc}
\usepackage[utf8]{inputenc}
\usepackage[polish]{babel}
\usepackage{lmodern}
\selectlanguage{polish}
\usepackage{graphicx}
\usepackage{biblatex}
\usepackage{csquotes}
\usepackage{booktabs}
\usepackage{tabularx}
\usepackage{xcolor}
\usepackage[hidelinks]{hyperref}
\addbibresource{bib.bib}
\setlength{\emergencystretch}{1.5em} % mniej wystających linii
% Łamanie długich adresów URL w bibliografii
\setcounter{biburllcpenalty}{7000}
\setcounter{biburlucpenalty}{8000}

\begin{document}

\pagenumbering{gobble}
\clearpage
\begin{figure}[h]
\centering
\includegraphics{media/ps-logo.png}
\end{figure}
\hspace{3cm}
\begin{center}Sprawozdanie z modułu nr 1\end{center}
\begin{center}CYBER 2026/2027\end{center}
\hspace{3cm}
\begin{center}\large\textbf{Cyberbezpieczeństwo}\end{center}
\hspace{7cm}
\begin{center}Credential stuffing -- przejmowanie kont
z użyciem danych uwierzytelniających z wycieków\end{center}
\begin{flushright}Kierunek: Informatyka
\end{flushright}
\begin{flushright}Członkowie zespołu:
\par
\textit{Dawid Pągowski}
\par
\textit{Bartłomiej Łokaj}
\par
\textit{Paweł Świderski}
\end{flushright}
\vfill
\begin{center}Gliwice, 2026/2027\end{center}

\newpage
\pagenumbering{arabic}
\tableofcontents

\newpage
\section{Wprowadzenie}

\subsection{Założenia dot. zadania laboratoryjnego}
Celem laboratorium jest analiza cyberzagrożenia credential stuffingu, czyli wykorzystywania wykradzionych poświadczeń (np. z wycieków) do nieautoryzowanego logowania na konta ofiar. Zagrożenie to zalicza się do kategorii \textit{przejęcie konta}.

\subsection{Uzasadnienie wyboru}

\begin{itemize}
    \item skala -- dodatkowe badania do raportu Verizon DBIR 2025 wskazują, że w logach dostawców usług SSO w medianie 19\% dziennych prób uwierzytelnienia wynikało z credential stuffingu. Ponadto w medianie tylko 49\% zapisanych haseł użytkownika było unikalnych, a pozostałe powtarzały się w różnych serwisach -- to właśnie wykorzystuje credential stuffing~\cite{verizon2025}.
    \item dobrze udokumentowane przypadki -- cyberzagrożenie jest opisane przez MITRE i OWASP~\cite{mitre_t1110004,owasp_cs}.
\end{itemize}

\subsection{Zakres pracy}

Sprawozdanie obejmuje: definicję zagrożenia i jego umiejscowienie w~podstawowych pojęciach bezpieczeństwa, opis mechanizmu ataku wraz z odniesieniem do macierzy MITRE ATT\&CK, analizę rzeczywistego incydentu (23andMe, 2023), ocenę skutków w~modelu CIA, metody wykrywania oraz rekomendowane zabezpieczenia. Praca opiera się wyłącznie na publicznie dostępnych źródłach; nie przeprowadzaliśmy żadnych testów na cudzych systemach.

\newpage

\section{Rozwinięcie}

\subsection{Definicja}

OWASP definiuje credential stuffing jako zautomatyzowane wstrzykiwanie skradzionych poświadczeń (par nazw użytkowników i haseł) do formularzy logowania w celu uzyskania nieautoryzowanego dostępu do kont~\cite{owasp_cs}.

Wyraźnie określa też credential stuffing jako podzbiór kategorii ataków brute force~\cite{owasp_cs}. Klasyczny atak brute force polega na zgadywaniu kolejnych kombinacji, a credential stuffing polega na próbowaniu \textbf{wyłącznie} znanych poświadczeń.

\subsection{Zagrożenie, podatność, incydent i ryzyko}

\begin{table}[h]
\centering
\small
\begin{tabularx}{\linewidth}{@{}l>{\raggedright\arraybackslash}X@{}}
\toprule
\textbf{Pojęcie} & \textbf{Odniesienie do credential stuffingu} \\
\midrule
Zagrożenie & Zautomatyzowane próby logowania danymi z wycieków, prowadzone na dużą skalę. \\
Podatność & Używanie tych samych haseł w wielu serwisach, brak uwierzytelniania wieloskładnikowego, brak ograniczeń liczby prób logowania. Jest to słabość procesu i konfiguracji, a nie błąd w kodzie. \\
Incydent & Udane zalogowanie się atakującego na konto ofiary, a w konsekwencji dostęp do jej danych lub wykonywanie operacji w jej imieniu. \\
Ryzyko & Wysokie prawdopodobieństwo (tanie i łatwo dostępne listy poświadczeń, automatyzacja) połączone ze skutkami zależnymi od wartości danych dostępnych z poziomu konta. \\
\bottomrule
\end{tabularx}
\caption{Credential stuffing w ujęciu podstawowych pojęć bezpieczeństwa}
\end{table}

\subsection{Mechanizm}

\begin{enumerate}
    \item Atakujący pozyskuje poświadczenia z dowolnego źródła: publiczne i niepubliczne wycieki, kampanie phishingowe, złośliwe oprogramowanie typu infostealer itd.~\cite{owasp_cs,hns_okta2024}.
    \item Atakujący wykorzystuje skradzione dane do masowych, zautomatyzowanych prób logowania na wielu różnych serwisach~\cite{owasp_cs}.
    \item Aby uniknąć blokad opartych na adresie IP, atakujący rozprasza ruch, np. przez sieć TOR lub proxy rezydencjalne, dzięki którym żądania wyglądają jak ruch zwykłych użytkowników domowych~\cite{hns_okta2024}.
    \item Pomyślne zalogowanie pozwala atakującemu wykorzystać konto i namierzyć kolejne ofiary lub okazje do ataków.
\end{enumerate}

Oprócz logowania się za pomocą poświadczeń, atakujący może również wykorzystać dane potrzebne do resetowania haseł: np. odpowiedzi na pytania zabezpieczające, numery telefonów, dane osobowe i inne~\cite{owasp_cs}.

\subsection{Potencjalne cele}
Cele można rozróżnić na użytkowników (bezpośrednie ofiary) i typy usług. W grupie użytkowników są to osoby korzystające z tych samych loginów i haseł w wielu miejscach. Do potencjalnych celów-usług należą wszystkie serwisy i aplikacje, które udostępniają publiczne formularze logowania~\cite{mitre_t1110004}.

\subsection{Porównanie z~pokrewnymi atakami}

\begin{table}[h]
\centering
\small
\begin{tabularx}{\linewidth}{@{}l>{\raggedright\arraybackslash}X>{\raggedright\arraybackslash}X@{}}
\toprule
\textbf{Atak} & \textbf{Co jest testowane} & \textbf{Źródło haseł} \\
\midrule
Brute force & Wiele haseł dla jednego konta & Słownik lub generowanie kombinacji \\
Password spraying & Jedno popularne hasło dla wielu kont & Listy najczęstszych haseł \\
Credential stuffing & Konkretne pary login--hasło & Wycieki z innych serwisów \\
\bottomrule
\end{tabularx}
\caption{Credential stuffing na tle pokrewnych ataków (opracowanie na podstawie~\cite{owasp_cheatsheet})}
\end{table}

\subsection{Macierz MITRE ATT\&CK}
\begin{table}[h]
\centering
\small
\begin{tabularx}{\linewidth}{@{}>{\raggedright\arraybackslash}p{3.2cm}>{\raggedright\arraybackslash}p{5.6cm}>{\raggedright\arraybackslash}X@{}}
\toprule
\textbf{Taktyka} & \textbf{Technika} & \textbf{Etap ataku} \\
\midrule
Reconnaissance & T1589.001 Gather Victim Identity Information: Credentials &
  Pozyskanie poświadczeń (etap 1) \\
Credential Access & T1110.004 Brute Force: Credential Stuffing &
  Masowe próby logowania (etap 2) \\
Command and Control & T1090.002 External Proxy, T1090.003 Multi-hop Proxy &
  Ukrywanie ruchu przez proxy rezydencjalne i TOR (etap 3) \\
Initial Access / Persistence & T1078 Valid Accounts &
  Korzystanie z~przejętych kont (etap 4) \\
\bottomrule
\end{tabularx}
\caption{Mapowanie etapów ataku na taktyki i~techniki MITRE ATT\&CK (opracowanie na podstawie~\cite{mitre_t1110004})}
\end{table}

Centralna jest technika T1110.004 z taktyki \emph{Credential Access}. Pozostałe pozycje opisują działania przed atakiem i po nim, co pokazuje, że credential stuffing jest zwykle jednym z etapów dłuższego łańcucha, a nie celem samym w sobie.

\subsection{Scenariusz ataku i przykład: 23andMe (2023)}

W październiku 2023 roku firma 23andMe, oferująca komercyjne badania genetyczne, padła ofiarą głośnego incydentu bezpieczeństwa, w wyniku którego doszło do masowego wycieku danych użytkowników. Przebieg tego ataku przebiegał według następującego schematu \cite{wiki23, sec23}:

\begin{itemize}
    \item \textbf{Zasób:} Baza kont użytkowników platformy oraz powiązana z nią funkcja społecznościowa \textit{DNA Relatives}, umożliwiająca odnajdywanie i łączenie krewnych na podstawie podobieństwa genetycznego \cite{wiki23, sec23}.
    \item \textbf{Sposób wejścia:} Atakujący wykorzystali technikę \textit{credential stuffing}, polegającą na automatycznym testowaniu na platformie 23andMe loginów i haseł, które wyciekły wcześniej z innych serwisów internetowych \cite{owasp, wiki23}.
    \item \textbf{Działanie atakującego:} Po uzyskaniu nieautoryzowanego dostępu do około 14 tysięcy kont (których właściciele stosowali te same hasła w różnych usługach), sprawcy wykorzystali funkcję \textit{DNA Relatives}. Pozwoliło im to na masowe zautomatyzowane pobieranie (tzw. \textit{scraping}) danych osobowych i genealogicznych wszystkich osób połączonych z przejętymi profilami. Działania te były ukierunkowane m.in. na użytkowników pochodzenia żydowskiego (aszkenazyjskiego) oraz chińskiego \cite{wiki23, sec23}.
    \item \textbf{Skutek:} Wyciek objął łącznie dane około \textbf{6,9 miliona} użytkowników. Ujawnione informacje zawierały m.in. imiona i nazwiska, rok urodzenia, płeć, dane o pochodzeniu etnicznym, drzewa genealogiczne oraz stopień pokrewieństwa \cite{sec23, wiki23}.
    \item \textbf{Wykrycie:} Incydent został wykryty w październiku 2023 roku, kiedy próbki skradzionych baz danych zaczęły być publikowane i oferowane na sprzedaż na forach hakerskich \cite{wiki23}.
    \item \textbf{Reakcja:} W odpowiedzi na incydent firma wymusiła reset haseł dla wszystkich użytkowników oraz wprowadziła obowiązkowe uwierzytelnianie dwuskładnikowe (2FA), aby zablokować możliwość dalszego wykorzystywania skradzionych poświadczeń \cite{wiki23}.
\end{itemize}

\subsection{Skutki}

Analiza incydentu w oparciu o triadę bezpieczeństwa informacji (model CIA) pokazuje selektywny charakter naruszenia \cite{wiki23, sec23}:

\begin{table}[h]
    \centering
    \begin{tabular}{p{3cm} p{2.5cm} p{8.5cm}}
        \toprule
        \textbf{Właściwość} & \textbf{Naruszona?} & \textbf{Uzasadnienie} \\
        \midrule
        Poufność & \textbf{Tak} & Doszło do nieuprawnionego ujawnienia wrażliwych danych osobowych, demograficznych i genetycznych 6,9 mln osób \cite{sec23}. \\
        Integralność & \textbf{Nie} & Dane w bazach 23andMe nie zostały zmodyfikowane, sfałszowane ani usunięte – atakujący jedynie je skopiowali \cite{wiki23}. \\
        Dostępność & \textbf{Nie} & Serwis i jego usługi pozostały w pełni dostępne dla użytkowników przez cały czas trwania incydentu (nie był to atak typu DoS ani ransomware) \cite{sec23}. \\
        \bottomrule
    \end{tabular}
    \caption{Ocena skutków wycieku danych 23andMe w modelu CIA.}
\end{table}

Poza bezpośrednim naruszeniem poufności, incydent wywołał szereg długofalowych konsekwencji dla przedsiębiorstwa:

\begin{itemize}
    \item \textbf{Skutki techniczne:} Konieczność pilnej przebudowy mechanizmów uwierzytelniania (wdrożenie obowiązkowego 2FA) oraz nałożenia limitów zapytań i dodatkowej kontroli dostępu do funkcji \textit{DNA Relatives}, aby zapobiec automatycznemu pobieraniu danych \cite{wiki23}.
    \item \textbf{Skutki organizacyjne:} Zaostrzenie wewnętrznych procedur reagowania na incydenty, przeprowadzenie zewnętrznych audytów bezpieczeństwa oraz powołanie specjalnych zespołów do obsługi roszczeń i komunikacji kryzysowej \cite{sec23}.
    \item \textbf{Skutki finansowe:} Drastyczny spadek wartości akcji spółki na giełdzie (o ponad 90\% względem szczytowych wycen), wysokie koszty obsługi prawnej i audytów, a także konieczność zawarcia ugody w ramach pozwów zbiorowych w USA na kwotę 30 milionów dolarów \cite{sec23}.
    \item \textbf{Skutki prawne:} Naruszenie przepisów o ochronie danych osobowych (w tym RODO oraz brytyjskiego UK GDPR), co skutkowało wszczęciem oficjalnych dochodzeń przez organy nadzorcze (m.in. brytyjskie ICO i kanadyjskie OPC) oraz nałożeniem kar finansowych za niewystarczające zabezpieczenie danych wrażliwych \cite{sec23}.
    \item \textbf{Skutki wizerunkowe:} Trwała utrata zaufania klientów. Ze względu na fakt, że danych genetycznych (w przeciwieństwie do hasła czy numeru karty płatniczej) nie da się zmienić, kryzys reputacyjny doprowadził do masowego usuwania kont przez użytkowników i załamania modelu biznesowego firmy \cite{wiki23}.
\end{itemize}

\subsection{Sposoby wykrywania}
TODO (osoba 3):
- oznaki w logach: duza liczba nieudanych zalgoowan na wiele kont z wiel uadreow iP
- ruch z seici anonimizujacych (np TOR)
- anomalie w zachowaniu klietna (MITRE T1110.004)

\subsection{Sposoby ograniczenia ryzyka}
TODO (osoba 3):
- zapobieganie (MFA, sprawdzanie hasel w bazach wycieków)
- reakcja: wymuszenie zmiany hasla, uniewaznienies sesji, powiadomienie userów

\newpage

\section{Podsumowanie i wnioski}

TODO (osoba 3):
- najwazniejsze ustalenia
- ocena skutecznosci zabezpieczen
- ktory etap najlatwiej przerwac
- ktore zabezpieczenie daje najwiekszy efekt
- jak sprawdzic czy ochrona dziala

\newpage
\section{Spis literatury}

\printbibliography[heading=none]

\end{document}
